\section{Listening and Adjusting}
% \begin{equation}
%     \dt{x} = x^2 \\
%     \frac{d}{dt} = x \\
%     dd=5
% \end{equation}
\subsection{The Kuramoto Rule}
A widely used model for interacting phases is
\begin{equation}
    \dt{\theta_{i}} = \omega_{i} + \frac{K}{N} \sum_{j=1}^{N} \sin(\theta_{j} - \theta_{i}), \quad i = 1, \dots , N
    \label{eqn:3}
\end{equation}
The three parts of eq.~\ref{eqn:3} have direct meanings:
\begin{description}[style=nextline]
    \item[Internal clock] $\omega_{i}$ 
is the rate the firefly would follow alone.
    \item[Coupling strength] $K\ge 0$ measures how strongly flashes influence one another
    \item[Timing correction] $\sin(\theta_{j} - \theta_{i})$ tells firefly $i$ whether another phase is ahead or behind
\end{description}

% \begin{itemize}
%     \item shuv
%     \item anni
%     \item sam
% \end{itemize}



% \begin{tasks}(2) %[label=\Roman*.](2)
%     \task Apple
%     \task Orange
%     \task Mango
%     \task Banana
% \end{tasks}

% \begin{tcolorbox}[title=Advantages,colback=blue!5,colframe=blue]
% \begin{itemize}
%     \item Fast
%     \item Portable
%     \item Open Source
% \end{itemize}
% \end{tcolorbox}

Table~\ref{tab:one} explains the correction term. The table is simple, but it contains the central intuition of the
model.

\begin{table}[h]
    \caption{How a neighboring firefly affects firefly $i$.}
    \label{tab:one}
    \centering
    \resizebox{\linewidth}{!}{%
    \begin{tabular}{cccl}
        \toprule
         \textbf{Relative pOsition} & \textbf{Phase Difference} & \textbf{Sine value} & \textbf{Effect on firefly $i$} \\
         \midrule
         Nrighbour ahead & $0<\theta_{j} - \theta_{i}<\pi$ & \cellcolor{lime}{\textbf{Positive}} & \cellcolor{lime}{Speeds up its phase} \\
         Phases aligned & $\theta_{j} - \theta_{i}=0$ & \cellcolor{blue!10}{\textbf{zero}} & \cellcolor{blue!10}{Makes no correcition} \\
         Neighbour behind & $-\pi<\theta_{j} - \theta_{i}<0$ & \cellcolor{red!10}{\textbf{Negative}} & \cellcolor{red!10}{Slows down its phase} \\
         Opoosite phases & $\theta_{j} - \theta_{i}=0$ & \cellcolor{blue!10}{\textbf{Zero}} & \cellcolor{blue!10}{Produces no preferred direction of correction.} \\
         \bottomrule 
    \end{tabular}
    }
\end{table}

The sine function is useful because it is smooth, periodic, and changes sign when the order of two phases
is reversed.

